\section{The Verbal Reward Stack}
\label{sec:rewardstack}

The learning-signal pillar (\S\ref{sec:learning}) examined how verbal feedback is compressed into parameter updates. This section asks the prior question: where does the reward itself come from? Training-time VRL has converged on four reward constructions, and we argue they are best read as a single progression: verifiable outcome rewards, rubric rewards, generative reward models, and process reward models (PRMs). Two properties order the rungs. The first is how much linguistic structure participates in computing the reward. The second, which we call the \emph{scalarization boundary}, is the point in that computation at which language collapses into the scalar the optimizer consumes. Ascending the stack moves the boundary later in the pipeline and attaches the signal at finer granularity, at rising verification cost. Table~\ref{tab:rewardstack} summarizes the progression; \S\ref{subsec:scalarization} makes the boundary explicit and traces its consequences.

\begin{table}
\caption{The verbal reward stack. Each rung is characterized by where linguistic structure collapses into the scalar consumed by the optimizer (the scalarization boundary) and by the granularity at which the signal attaches to the trajectory.}
\label{tab:rewardstack}
\small
\begin{tabular}{@{}p{0.14\textwidth}p{0.27\textwidth}p{0.22\textwidth}p{0.1\textwidth}p{0.15\textwidth}@{}}
\toprule
Rung & Reward computation & Scalarization boundary & Granularity & Representative work \\
\midrule
Verifiable outcome & Programmatic check against ground truth & At the environment; no linguistic structure in the reward & Episode & \citep{lambert2024tulu,guo2025deepseek} \\
Rubric reward & Judge checks instance-specific criteria; item scores pooled & At aggregation of criterion verdicts & Response & \citep{gunjal2025rubrics,viswanathan2025checklists,huang2025reinforcement} \\
Generative reward model & Judge writes principles and critique, then a verdict & At verdict emission; critique retained as an audit trail & Response & \citep{zhang2025genrm,liu2025inferencetime,chen2025rmr1} \\
Process reward & Step-level verification, increasingly verbalized & Per step (per token at the limit) & Step or token & \citep{wang2024mathshepherd,khalifa2025processrewardmodelsthink,xu2026rubrics} \\
\bottomrule
\end{tabular}
\end{table}

\subsection{Verifiable Outcome Rewards}
\label{subsec:rlvr}

The bottom rung is reinforcement learning with verifiable rewards (RLVR), a term introduced by the Tulu 3 post-training recipe, which combined supervised fine-tuning, preference optimization, and RL against programmatic correctness checks in a fully open pipeline \citep{lambert2024tulu}. DeepSeek-R1 supplied the flagship result: pure RL against verifiable outcomes, with no human-labeled reasoning trajectories, produced emergent self-reflection and verification behavior and outperformed counterparts trained by supervised learning on human demonstrations across verifiable mathematics, coding, and STEM tasks \citep{guo2025deepseek}. In the vocabulary of \S\ref{sec:framework}, RLVR is the degenerate case of the stack: scalarization happens at the environment itself, and no linguistic structure enters the reward. We include it because it is the baseline every higher rung generalizes, and because its limitations motivated the climb. Verifiable rewards are sparse in agentic settings; Agent-RLVR densifies them with teacher-style verbal guidance, plans and error feedback attached to failed attempts, lifting pass@1 on SWE-Bench Verified from 9.4\% to 22.4\%, and to 27.8\% with a reward model trained on the guided data \citep{da2025agentrlvr}. Verifiable rewards also reach only domains with checkable answers, the gap the rubric rung addresses. Even within their domain, measured gains deserve caution: a position paper on RLVR's hidden costs shows that budget mismatches, attempt inflation with calibration drift, and benchmark contamination confound many headline results, several of which shrink substantially or disappear under budget-matched reproduction \citep{wu2025position}. We revisit this measurement problem in \S\ref{sec:evaluation}; it matters here because every rung above inherits RLVR's optimization machinery, and often its evaluation habits.

\subsection{Rubric Rewards}
\label{subsec:rubrics}

Rubric rewards extend verifiability to open-ended domains by replacing the ground-truth check with instance-specific criteria that a judge evaluates and an aggregation step pools into a scalar. Rubrics as Rewards is the anchor method: scoring rollouts against per-instance rubrics yields relative gains of up to 31\% on HealthBench and 7\% on GPQA-Diamond over LLM-as-judge baselines that assign direct Likert scores, and the structured criteria let smaller judges score more consistently \citep{gunjal2025rubrics}. Checklist feedback makes the same move for instruction following: per-instruction criteria scored by AI judges and verifier programs form an RL reward that is the only method in its comparison to improve on every benchmark, including a four-point gain in hard satisfaction rate on FollowBench \citep{viswanathan2025checklists}.

Rubric supply creates a three-way choice between coverage, adaptation, and verification cost. Large static banks dominate when the task distribution is broad but stable: human and synthetic collections scale open-ended RL, while coarse-to-fine or contrastive generation reduces authoring cost and derives criteria from examples \citep{huang2025reinforcement,li2026rubrichub,liu2025openrubrics}. Dynamic rubrics dominate when the policy or case mix changes during training, because fixed criteria become stale and expose exploitable omissions. Online generation, policy-coupled evolution, self-rewarding, and case-conditioned criteria all trade additional judge calls for criteria matched to the current failure distribution \citep{rezaei2025online,guan2026evorubric,ye2025selfrewarding,wang2025infimedorbit}. Selective verification is preferable when checking every criterion costs more than generating candidates: an adversarial critic can surface only likely violations for external validation and outperform exhaustive checking with fewer verifications \citep{wu2025rlac}. The same logic applies even in domains with correct answers. Process-oriented rubrics catch invalid derivations that outcome checks miss, raising Verified Pass@1024 on AIME2024 from 26.7\% to 62.6\% while reducing ``Miracle Steps'' by 71\% \citep{yuan2025curing}. Static banks therefore buy amortized coverage, dynamic rubrics buy robustness to distribution shift, and selective or process checks buy precision where verification is expensive or outcome correctness is insufficient.

The closest prior survey treats rubrics as a framework recurring across evaluation, RL, and safety, organized into evaluative, training, and intrinsic levels \citep{chen2026from}. Our reading is complementary but different: we place rubrics as one rung of a reward stack ordered by scalarization and granularity, which exposes what they share with the verifiers below and the generative judges above. For this rung, scalarization occurs at aggregation: linguistic structure survives through criterion checking, but once item verdicts are pooled, anything the rubric failed to anticipate is invisible to the training signal, a slack we return to in \S\ref{subsec:scalarization}.

\subsection{Generative Reward Models}
\label{subsec:genrm}

The third rung verbalizes the reward computation itself. Generative verifiers recast reward modeling as next-token prediction: the verifier writes a verification chain of thought before its verdict, inheriting instruction following and majority voting over sampled rationales, and improves best-of-N accuracy from 5\% to 45.3\% on algorithmic tasks and from 73\% to 93.4\% on GSM8K \citep{zhang2025genrm}. DeepSeek-GRM extends this to generalist reward modeling through self-principled critique tuning, online RL that teaches the model to generate principles and then critiques, and shows that sampling more reward computations at inference, guided by a meta reward model, can outperform scaling the reward model's size at training time \citep{liu2025inferencetime}. A growing cluster treats verdict production as a reasoning task in its own right: chain-of-rubrics reward models self-generate rubrics or reference solutions before judging, outperforming far larger open and proprietary models by up to 4.9\% \citep{chen2025rmr1}; reward reasoning models spend adaptive test-time compute on hard judgments \citep{guo2025reward}; J1 converts judgment itself into a verifiable-reward RL problem \citep{whitehouse2025j1}; EvalPlanner decouples evaluation planning from execution, reaching 93.9 on RewardBench \citep{saha2025learning}; and further variants optimize judgment-thinking traces with offline and online RL \citep{huang2025thinkj}, extend judges to long-horizon reasoning \citep{hong2025thinkrm}, or pair a reasoning reward model with a dedicated benchmark (\S\ref{sec:evaluation}) \citep{zhou2025libra}.

Supervision for the judges is increasingly self-generated. Self-taught evaluators iterate without any human preference labels, improving a 70B judge from 75.4 to 88.3 on RewardBench \citep{wang2024selftaught}; jointly predicting critiques and scalar rewards improves reward accuracy by 3.7\%--7.3\% \citep{yu2024selfgenerated}; self-training on unlabeled data yields foundation reward models that emit labels with rationales \citep{wang2025gramr2}; and meta-rewarding has the model judge its own judgments, the recursive rung of the stack, lifting the AlpacaEval~2 win rate from 22.9\% to 39.4\% \citep{wu2024metarewarding}. The critique, however, is not automatically load-bearing: generative reward models trained only on verdict correctness learn to guess right outcomes from unsound critiques, and supervising the critique itself, by rewarding similarity to human critiques, outperforms outcome-only training \citep{wang2026reward}, closing the loop between the reward layer and the human verbal feedback of \S\ref{sec:learning}. For this rung, scalarization occurs at verdict emission: the verbal computation demonstrably improves the verdict and leaves an auditable trail, but the optimizer still consumes a scalar unless the critique is separately routed into revision or a language-space update (\S\ref{sec:optimizers}).

\subsection{Process Reward Models}
\label{subsec:prm}

The fourth rung moves the signal inside the trajectory. PRMs descend from human step-level annotation \citep{lightman2024let}, which automation quickly replaced: Math-Shepherd constructs step labels automatically and lifts GSM8K accuracy from 77.9\% to 84.1\% through step-by-step PPO \citep{wang2024mathshepherd}, and OmegaPRM's divide-and-conquer tree search locates the first error by binary search, collecting over 1.5 million step annotations and improving MATH500 from 51\% to 69.4\% with weighted self-consistency \citep{luo2024improve}. The subsequent diagnosis was sobering. ProcessBench, an error-localization benchmark whose protocol \S\ref{sec:evaluation} describes, finds that existing PRMs typically fail to generalize beyond GSM8K/MATH difficulty and underperform prompted critic models \citep{zheng2024processbench}. A companion study explains why: Monte Carlo-estimated step labels underperform LLM-as-a-judge and human annotation because completion models verify steps inaccurately, and best-of-N evaluation quietly drifts PRMs back toward outcome scoring; consensus filtering between the two estimators significantly improves performance and data efficiency \citep{zhang2025lessons}.

Generative PRMs are where the judge line and the step line converge. ThinkPRM verifies every step by generating a verification chain of thought, outperforming discriminative verifiers with only 1\% of the PRM800K labels \citep{khalifa2025processrewardmodelsthink}; StepWiser reframes stepwise reward modeling from classification to reasoning, training the judge with RL from rollout outcomes \citep{xiong2025stepwiser}; and DG-PRM maintains a tree of fine-grained reward criteria selected dynamically per step \citep{yin2025dynamic}. SPARK closes the loop for domains without verifiable answers: synthetic step-level verification data trains a PRM that then drives RL, beating ground-truth RLVR at 47.4\% versus 43.9\% averaged over six mathematics benchmarks \citep{rahman2025spark}, and OmegaPRM-style pipelines now extend to multimodal reasoning with over 700k automatically generated step annotations \citep{du2025mmprm}. At the finest granularity, response-level rubrics can be converted into token-level credit through a relevance discriminator \citep{xu2026rubrics}, at which point the reward stack meets the credit-assignment problem of \S\ref{sec:credit}.

\subsection{The Scalarization Boundary}
\label{subsec:scalarization}

We can now state the section's organizing claim plainly. Every rung ends the same way: a scalar crosses the interface to the optimizer. What distinguishes the rungs is where the collapse happens. Verifiers scalarize at the environment; rubrics at aggregation, after criterion-level checking; generative reward models at verdict emission, after a full verbal computation; PRMs per step; token-level methods per token. What is lost at the boundary is everything the verbal computation established that a magnitude cannot carry: which failure occurred, what would repair it, and why the trajectory earned its score. The gradient sees that a rollout scored 0.3, not that it scored 0.3 because its third step asserted an unsupported lemma. Ascending the stack narrows the loss by shrinking the span each scalar summarizes, but only the language-space optimizers of \S\ref{sec:optimizers} remove the boundary altogether, applying feedback as text without collapsing it. In the terms of \S\ref{sec:framework}, rubric scores and PRM verdicts are verbal reward functions, language occupying the role of the reward function in the underlying decision process; the practitioner question of which rung to adopt, a tradeoff between signal granularity and verification cost, is taken up in \S\ref{sec:practitioner}.

The boundary is also where reward hacking concentrates, because whatever the verbal computation understood but the scalar cannot transmit becomes exploitable slack. The evidence spans every rung. Outcome rewards breed Miracle Steps \citep{yuan2025curing}, and static rubrics are gameable by design \citep{rezaei2025online}. A systematic study of rubric-based RL locates the exploitation in precisely this slack: when rubrics leave failure modes unspecified, rubric-based verifiers prefer the trained checkpoint while rubric-free judges prefer the base model \citep{mahmoud2026reward}, a finding set \S\ref{sec:safety} examines in full. One rung up, generative reward models guess correct verdicts from unsound critiques \citep{wang2026reward}; one rung further, mis-estimated step labels corrupt PRM training \citep{zhang2025lessons}; and at the base, headline RLVR gains partly dissolve under budget-matched evaluation \citep{wu2025position}. The emerging defenses sit, tellingly, at the boundary itself: cross-family judge panels with a verifier-free self-internalization diagnostic \citep{mahmoud2026reward}, meta-verification and peer consensus over self-generated rubrics \citep{guan2026evorubric}, format constraints on PRM-driven RL \citep{rahman2025spark}, and consensus filtering of step labels \citep{zhang2025lessons}. Whether these patches suffice, or whether the boundary itself must move, is the safety question of \S\ref{sec:safety}.
